\section{Introduction}

Lung adenocarcinoma (LUAD) develops through a stereotypic morphological sequence that
begins with atypical adenomatous hyperplasia (AAH) and proceeds through adenocarcinoma
\emph{in situ} (AIS) and minimally invasive adenocarcinoma (MIA) to invasive
adenocarcinoma (IAC). Each stage is histologically defined, and the transition from
precursor lesion to invasive carcinoma is the point at which clinical management changes.
The sequence therefore provides a well-defined system in which to ask what transcriptional
changes accompany invasion and in what order they arise.

Most transcriptomic descriptions of this sequence derive either from bulk tissue
\citep{tcga2014luad} or from single-cell profiling restricted to the invasive tumour and
adjacent normal lung, and both designs conflate the intermediate states that define the
sequence. Single-cell atlases spanning the full spectrum are available \citep{han2024}, but
these rarely carry a spatially resolved counterpart, so the tissue architecture surrounding
a given transcriptional state, such as the niche composition around an alveolar type II
cell, remains unobserved.

Two obstacles are specific to this disease. Tumour cells are difficult to identify without
copy-number information, because the reference panels used for spatial deconvolution
contain no malignant cell type; malignant epithelium is assigned to the most similar normal
epithelial class, leaving any composition-based description of a niche structurally blind
to the tumour. In addition, invasive regions of a section are both more cellular and more
transcriptionally active than precursor regions, so comparisons between the two are also
comparisons of RNA content.

We assembled paired single-nucleus RNA sequencing and Visium spatial transcriptomics from
23 patients spanning all five stages. The single-cell data define cell states and a
patient-internal disease axis. The spatial data define tissue niches
\citep{singhal2024banksy,cable2022rctd} and provide the context in which the copy-number
signal is placed \citep{cabrejas2026fastcnv}. The same disease axis is then used as a query
against two perturbation libraries \citep{subramanian2017}, to ask whether any known
compound or gene perturbation reverses it.

A methodological result obtained early shaped the remainder of the study. Sequencing depth
and tissue density are not separable in this cohort, and their coupling is strong enough to
generate stage-dependent patterns where none exist. Controlling for it, and reporting what
survives when that control is incomplete, became the organising constraint of the work. We
therefore present positive and negative results together and record the analytical
difficulties encountered, several of which required the withdrawal of an apparently sound
result.
